\subsection{Two-Stage Position Estimator}
\label{subsec:estimator}
The separable structure~\eqref{eq:separable} suggests estimating the direction and the distance in two stages from the same $K$ measurements. The direction and the common amplitude are obtained by weighted nonlinear least squares (NLS):
\begin{equation}
 (\hat{\vect u},\hat\beta)=\arg\min_{\vect u\in\mathbb S^2,\ \beta>0}\ \sum_{i=1}^{K}w_i\left(y_i-\beta R(\vect n_i\trans\vect u)\right)^2,
 \label{eq:nls}
\end{equation}
with $w_i=L/\sigma_i^2$, equal for all orientations in the experiments. For a fixed $\vect u$ the amplitude has the closed form $\hat\beta(\vect u)=\sum_iw_iy_iR_i/\sum_iw_iR_i^2$, $R_i=R(\vect n_i\trans\vect u)$, so~\eqref{eq:nls} reduces to a search over the unit sphere. It is solved by a Levenberg--Marquardt iteration on the sphere, started from the best three of a fixed set of seed directions; solutions with competing minima are flagged (Appendix~\ref{app:estimators}). This stage uses only the shape of the measurements: it does not depend on $A$, on $m$ or on the LED pattern. As alternatives, the ratios to the strongest orientation, $y_i/y_j=R(\vect n_i\trans\vect u)/R(\vect n_j\trans\vect u)$, which do not depend on $\beta$, can be fitted by generalized (GLS) or weighted (WLS) least squares, with the full or the diagonal first-order covariance of the ratios. For a cosine-power response, the roots of the ratios are linear in $\vect u$ and the GLS and WLS directions are eigenvectors, as derived in~\cite{Acuna2026} for a reoriented transmitter. The SQ response admits no such linearization, and its ratios are fitted iteratively on the sphere. Appendix~\ref{app:estimators} gives the estimators for an arbitrary response, as used here with SQ, together with Algorithm~\ref{alg:estimator}. \add{All of them accept any set of known normals: they require neither a cone nor uniformly spaced or exactly commanded orientations, and they use the normals actually measured.}

The distance then follows in closed form from the amplitude:
\begin{equation}
 \hat d=\sqrt{\frac{A\cos^{m}\hat\phi}{\hat\beta}},\qquad
 \cos\hat\phi=-\vect n_t\trans\hat{\vect u},\qquad
 \hat{\vect r}=\vect t-\hat d\,\hat{\vect u}.
 \label{eq:range}
\end{equation}
Given $A$ and $m$, the map $(\vect u,\beta)\mapsto(\vect u,d)$ is one-to-one, so the two stages yield the maximum-likelihood position under~\eqref{eq:observation}. Unlike the transmitter-steering system of~\cite{Acuna2026}, which obtained the range from an additional measurement with the LED and PD aligned, no cooperative measurement is required: direction and range come from the same reorientation sequence. \add{The estimator is deterministic and model-based: besides the $K$ RSS and the normals, it uses only the calibrated parameters of Table~\ref{tab:parameters}, without a fingerprint database, a training stage or a stochastic search, and it takes a few milliseconds per position (Section~\ref{subsec:res-time}).}

The range inherits the errors of the direction through the LED pattern. Differentiating~\eqref{eq:range} with respect to the irradiance angle gives:
\begin{equation}
 \frac{\partial\ln d}{\partial\phi}=-\frac{m}{2}\tan\phi,
 \label{eq:range-sensitivity}
\end{equation}
so that, with $m\approx6$ and $\phi=25^\circ$, a direction error of $1^\circ$ changes the distance by 2.4\,\%, about 4~cm at 1.8~m. Likewise, a mismatch between the calibrated and true values of $A$ or of the pattern scales the distance by $\sqrt{\hat A\cos^{\hat m}\phi/(A\cos^{m}\phi)}$, a bias that no amount of averaging removes. These two mechanisms, rather than the measurement noise, will turn out to limit the accuracy (Section~\ref{subsec:res-bound}).

\subsection{RSS-Shape Test}
\label{subsec:shape}
At a known position, the ratios $q_i=y_i/R(\vect n_i\trans\vect u)$ must all equal $\beta$. Their spread:
\begin{equation}
 \rho=100\sqrt{\frac1K\sum_{i=1}^{K}\left(\frac{q_i}{\bar q}-1\right)^2}\quad(\%),\qquad
 \bar q=\frac1K\sum_{i=1}^{K}q_i,
 \label{eq:shape}
\end{equation}
pooled over several positions, tests the angular part of the model independently of $A$, $m$ and the LED pattern, which are common to the $K$ orientations. It serves three purposes. With known positions, it discriminates between candidate receiver responses, and it calibrates \del{the angular zero $\theta_0$ of the rotation stage (and verifies the adapter tilt)}\add{the orientation parameters of the receiver, such as the angular zero of a rotation stage and the tilt of the PD (Section~\ref{subsec:hardware})}, by minimizing $\rho$ over \del{the orientation parameters}\add{them}. With unknown positions, the same statistic evaluated at the estimate, $\hat\rho$, is the residual of~\eqref{eq:nls}; it requires no ground truth and is used as an integrity indicator (Section~\ref{subsec:res-integrity}).

\subsection{LED Calibration From a Flat Scan}
\label{subsec:led-calibration}
With the PD facing up and scanning a horizontal plane at a vertical distance $h$ below the LED, the incidence and irradiance angles coincide, $\cos\psi=\cos\phi=h/d$, and the logarithm of~\eqref{eq:rss-model} is linear in the unknowns:
\begin{equation}
 \ln\frac{y\,d^2}{R(\phi)}=\ln A+m\ln\cos\phi.
 \label{eq:led-calibration}
\end{equation}
Given the LED position, $A$ and $m$ follow from linear least squares over the scan points (joint estimate), or from the point closest to the LED axis for $A$ and least squares for $m$ (staged estimate). The horizontal position of the LED is found by minimizing the residual of the joint fit. Its height, however, is not identifiable from a single plane: a higher LED is seen at smaller angles and longer distances, and larger values of $A$ and $m$ reproduce almost the same map. It must be measured, or calibrated with a second plane.

Because~\eqref{eq:led-calibration} is linear, the number and placement of calibration points can be studied exactly on measured data (Section~\ref{subsec:res-budget}). Points close to the LED axis determine $A$ well but carry little information on $m$; the calibration points must cover the emission angles at which the system operates.
